\subsection{Теорема о промежуточных значениях непрерывной на отрезке функции.}

\subsubsection*{Теорема Больцано-Коши (О промежуточном значении / О нулях)}
\textbf{Теорема.} Если $f(x) \in C([a, b])$ и на концах отрезка принимает значения разных знаков ($f(a) \cdot f(b) < 0$), то найдется точка $c \in (a, b)$, такая что $f(c) = 0$.

\begin{tcolorbox}[title=Доказательство (Метод деления отрезка пополам)]
	Положим $[a_0, b_0] := [a, b]$. По условию $f(a_0) \cdot f(b_0) < 0$.

	1. \textbf{Шаг построения.}
	Предположим, что мы уже построили отрезок $[a_k, b_k]$, такой что $f(a_k) \cdot f(b_k) < 0$.
	Найдем его середину $c_k = \frac{a_k + b_k}{2}$.
	\begin{itemize}
		\item Если $f(c_k) = 0$, то точка найдена, доказательство завершено.
		\item Если $f(c_k) \neq 0$, то рассматриваем произведение $f(a_k) \cdot f(c_k)$:
		\begin{enumerate}
			\item Если $f(a_k) \cdot f(c_k) < 0$ (знаки на концах левой половины разные), то выбираем $[a_{k+1}, b_{k+1}] = [a_k, c_k]$.
			\item Если $f(a_k) \cdot f(c_k) > 0$ (знаки одинаковые), то так как $f(a_k)$ и $f(b_k)$ имеют разные знаки, $f(c_k)$ и $f(b_k)$ обязаны иметь разные знаки. То есть $f(c_k) \cdot f(b_k) < 0$. Выбираем $[a_{k+1}, b_{k+1}] = [c_k, b_k]$.
		\end{enumerate}
	\end{itemize}

	2. \textbf{Система вложенных отрезков.}
	Если процесс не оборвался на шаге 1, мы получаем последовательность отрезков:
	\[ [a_0, b_0] \supset [a_1, b_1] \supset \dots \supset [a_k, b_k] \supset \dots \]
	Свойства этой последовательности:
	\begin{itemize}
		\item Длина стремится к нулю: $|b_k - a_k| = \frac{b-a}{2^k} \to 0$.
		\item На концах каждого отрезка значения функции имеют разные знаки:
		\[ \forall k \in \mathbb{N}: \quad f(a_k) \cdot f(b_k) < 0 \]
	\end{itemize}
	
	3. \textbf{Предельный переход.}
	По принципу вложенных отрезков (Лемма Кантора) существует единственная точка $c$, принадлежащая всем отрезкам:
	\[ c = \bigcap_{k=0}^{\infty} [a_k, b_k], \quad \text{причем } \lim_{k\to\infty} a_k = c, \;\; \lim_{k\to\infty} b_k = c \]
	
	В силу непрерывности функции $f(x)$:
	\[ \lim_{k\to\infty} f(a_k) = f(c) \quad \text{и} \quad \lim_{k\to\infty} f(b_k) = f(c) \]
	
	Перейдем к пределу в неравенстве $f(a_k) \cdot f(b_k) < 0$ (при переходе к пределу строгое неравенство может стать нестрогим):
	\[ \lim_{k\to\infty} (f(a_k) \cdot f(b_k)) \le 0 \implies f(c) \cdot f(c) \le 0 \implies (f(c))^2 \le 0 \]
	
	Так как квадрат действительного числа не может быть отрицательным, остается единственный вариант:
	\[ (f(c))^2 = 0 \implies f(c) = 0 \]
	Точка $c$ найдена.
\end{tcolorbox}

\subsubsection*{Следствие}

\textbf{Теорема.} Если $f(x) \in C[a, b]$ (непрерывна), то $f(x)$ на $[a, b]$ принимает все значения между $f(a)$ и $f(b)$ (то есть любое значение из отрезка с концами $f(a)$ и $f(b)$).

\begin{tcolorbox}[title=Доказательство]
    Пусть $C$ — число между $f(a)$ и $f(b)$.
    
    \begin{enumerate}
        \item Случай $f(a) = f(b)$. 
        Тогда промежуток вырождается в точку. Значение $C = f(a)$ принято на концах.
        
        \item Случай $f(a) \neq f(b)$.
        \begin{enumerate}
            \item[а)] Значения $f(a)$ и $f(b)$ (граничные) приняты в точках $a$ и $b$.
            
            \item[б)] Пусть $C$ находится \textbf{строго} между $f(a)$ и $f(b)$.
            Рассмотрим вспомогательную функцию:
            \[ g(x) = f(x) - C \]
            Так как $f(x) \in C[a, b]$, то и разность $g(x) \in C[a, b]$.
            
            Проверим знаки на концах (для определенности пусть $f(a) < C < f(b)$):
            \[ g(a) = f(a) - C < 0 \]
            \[ g(b) = f(b) - C > 0 \]
            
            Функция непрерывна и меняет знак на концах отрезка. 
            По теореме Больцано-Коши (о нулях непрерывной функции):
            \[ \exists c \in (a, b): \;\; g(c) = 0 \]
            \[ f(c) - C = 0 \implies f(c) = C \]
        \end{enumerate}
    \end{enumerate}
    \textbf{Вывод:} Любое промежуточное значение достигается в некоторой точке $c$.
\end{tcolorbox}




